\section{RESUMEN}

Este trabajo evaluó la presencia molecular y la caracterización genómica de SARS-CoV-2, poliovirus y virus sincitial respiratorio (RSV) en aguas servidas de la Planta de Tratamiento de Aguas Residuales (PTAR) Quitumbe, en Quito, con el fin de cubrir la limitada disponibilidad de información multipatógeno en aguas residuales del país. Las muestras se recolectaron mediante un dispositivo de muestreo pasivo tipo torpedo durante 56 semanas; se concentraron con polietilenglicol, se extrajo el ARN y se aplicó RT-qPCR para la detección, PCR multiplex con cebadores ARTIC para SARS-CoV-2 y RSV, y PCR anidada de la región VP1 para poliovirus, seguidas de secuenciación Oxford Nanopore y análisis bioinformático. Para SARS-CoV-2 se detectó material viral en la mayoría de las 42 muestras semanales, y seis muestras tempranas, con los valores de Cq más bajos, permitieron asignar linajes descendientes de JN.1 de los clados 24A y 24H. El esfuerzo de secuenciación posterior no rindió caracterización debido a la baja carga viral y a los valores de Cq tardíos. En las 22 muestras analizadas para poliovirus no se detectó el virus, pero se identificó Coxsackie A, un enterovirus no polio; en RSV se obtuvo señal molecular en cuatro muestras sin caracterización del subtipo. Las muestras retrospectivas conservadas en DNA/RNA Shield 2X y en ARN a −20~°C no rindieron secuencia. Los resultados confirman la factibilidad de un esquema de vigilancia multipatógeno en un sitio centinela urbano y señalan la carga viral y la conservación de las muestras como los factores que condicionan la caracterización genómica. Se recomienda el procesamiento temprano de las muestras, la selección por Cq previa a la secuenciación y la integración con datos clínicos para fortalecer la vigilancia ambiental en Quito.

\vspace{0.5\baselineskip}

\noindent\textbf{Palabras clave:} vigilancia basada en aguas residuales, SARS-CoV-2, poliovirus, virus sincitial respiratorio, secuenciación Oxford Nanopore, epidemiología genómica, enterovirus, PTAR Quitumbe.
